\documentclass[UTF8,11pt,a4paper]{ctexart}
\usepackage[a4paper,left=18mm,right=18mm,top=20mm,bottom=18mm]{geometry}
\usepackage{fontspec}
\usepackage{xcolor,graphicx,tabularx,booktabs,array,fancyhdr,hyperref,xparse,enumitem,float,amsmath}
\setmainfont{Arial}
\setCJKmainfont[Path=C:/Windows/Fonts/]{msyh.ttc}
\setCJKsansfont[Path=C:/Windows/Fonts/]{simhei.ttf}
\definecolor{navy}{HTML}{0B1730}
\definecolor{ink}{HTML}{15253B}
\definecolor{muted}{HTML}{66758A}
\definecolor{teal}{HTML}{0D9D91}
\definecolor{blue}{HTML}{2F68E8}
\definecolor{orange}{HTML}{E68A00}
\definecolor{violet}{HTML}{7655D7}
\definecolor{line}{HTML}{DCE5EF}
\definecolor{paleBlue}{HTML}{EEF4FF}
\definecolor{paleTeal}{HTML}{EAF9F6}
\definecolor{paleGold}{HTML}{FFF8E8}
\definecolor{paper}{HTML}{F7FAFC}
\setlength{\parindent}{0pt}
\setlength{\parskip}{4.5pt}
\setlength{\headheight}{23pt}
\setlist{nosep,leftmargin=1.5em}
\hypersetup{colorlinks=true,linkcolor=blue,urlcolor=blue,pdfauthor={Galen Research Workbench},pdftitle={07—其他材料：多智能体对比实验与系统测试报告}}
\pagestyle{fancy}
\fancyhf{}
\fancyhead[L]{\small\color{muted}\textbf{GALEN}  /  07—其他材料}
\fancyhead[R]{\small\color{muted}多智能体对比实验与系统测试报告}
\fancyfoot[L]{\small\color{muted}XH-202620  ·  评审证据材料}
\fancyfoot[R]{\small\color{muted}\thepage}
\renewcommand{\headrulewidth}{0.4pt}
\renewcommand{\footrulewidth}{0.4pt}
\makeatletter
\renewcommand{\headrule}{\hbox to\headwidth{\color{line}\leaders\hrule height \headrulewidth\hfill}}
\renewcommand{\footrule}{\hbox to\headwidth{\color{line}\leaders\hrule height \footrulewidth\hfill}}
\makeatother
\newsavebox{\cardbox}
\NewDocumentEnvironment{card}{m +b}{%
  \par\noindent\begin{lrbox}{\cardbox}\begin{minipage}{0.935\linewidth}#2\end{minipage}\end{lrbox}%
  \fcolorbox{line}{#1}{\usebox{\cardbox}}\par
}{}
\NewDocumentEnvironment{darkbox}{}{%
  \par\noindent\begin{lrbox}{\cardbox}\begin{minipage}{0.935\linewidth}\color{white}
}{\end{minipage}\end{lrbox}\fcolorbox{navy}{navy}{\usebox{\cardbox}}\par}
\newcommand{\kicker}[1]{\textcolor{teal}{\sffamily\bfseries\small #1}\par\vspace{1mm}}
\newcommand{\metric}[3]{\begin{minipage}[t]{0.30\linewidth}\centering{\color{#1}\fontsize{23}{27}\selectfont\sffamily\bfseries #2}\par{\small\color{muted}#3}\end{minipage}}
\newcommand{\statusok}{\textcolor{teal}{\bfseries 通过}}
\newcommand{\code}[1]{\texttt{\small #1}}
\newcolumntype{Y}{>{\raggedright\arraybackslash}X}
\begin{document}

\thispagestyle{empty}
\pagecolor{paper}
\vspace*{11mm}
{\color{teal}\sffamily\bfseries\small GALEN  /  XH-202620  /  2026.09}\par
\vspace{7mm}
{\color{navy}\fontsize{28}{35}\selectfont\sffamily\bfseries 07—其他材料}\par
{\color{navy}\fontsize{23}{31}\selectfont\sffamily\bfseries 多智能体对比实验与系统测试报告}\par
\vspace{4mm}
{\color{muted}\Large 用可重复任务、运行轨迹与交付物证明系统能力}\par
\vspace{11mm}
\begin{card}{white}
\vspace{2mm}
\metric{blue}{98.7\%}{Galen 硬门完成率}\hfill
\metric{teal}{2.11}{平均工具调用 / 次}\hfill
\metric{violet}{9.2 s}{中位端到端时延}
\vspace{2mm}
\end{card}
\vspace{6mm}
\begin{card}{paleTeal}
\textbf{本报告回答评委最关心的三个问题}\par
Galen 是否比通用 Agent 更稳定地完成康复科研任务？专用架构的优势来自哪里？结果能否回到任务卡、原始运行与正式产物逐项核验？
\end{card}
\vspace{5mm}
\begin{tabularx}{\textwidth}{@{}Y Y@{}}
\begin{card}{paleBlue}\textbf{主对比实验}\par 15 张任务卡 × 5 次重复 × 2 套完整系统，共 150 次对齐运行。\end{card}&
\begin{card}{paleGold}\textbf{架构消融与回归}\par 上下文压缩、无状态、执行契约、长任务看门狗和工程测试。\end{card}
\end{tabularx}
\vfill
\begin{card}{white}
\textbf{代码版本}\quad \code{galen-research-workbench / 3a179fb}\qquad
\textbf{公开仓库}\quad \href{https://github.com/Drehabwen/Galen}{github.com/Drehabwen/Galen}\par
\textbf{事实源}\quad 原始 JSONL、冻结任务卡、指标 CSV、统计脚本与 LaTeX 源文件
\end{card}
\clearpage
\pagecolor{white}

\kicker{01  /  结论先行}
\section*{不是“能回答”，而是“能按科研约束交付”}
\begin{card}{paleBlue}
\textbf{在本次冻结的康复科研任务集上，Galen 完成 74/75 次严格验收；Codex + DeepSeek 完成 21/75 次。}\par
两组使用相同的 15 张任务卡、固定输入、重复编号与 DeepSeek V4 Flash 模型标签。差异主要出现在项目范围保持、数据治理、工具预算与正式产物交付，而不是文案流畅度。
\end{card}
\begin{tabularx}{\textwidth}{@{}p{34mm}Y Y@{}}
\toprule
\textbf{观察维度} & \textbf{Galen} & \textbf{通用 CLI Agent 基线}\\
\midrule
硬门完成率 & 74/75，98.7\% & 21/75，28.0\%\\
工具经济性 & 平均 2.11 次/运行 & 平均 6.12 次/运行\\
端到端时延 & 中位 9.2 秒 & 中位 19.3 秒\\
上下文修订 & 结构化项目状态吸收修订 & 易重新带回旧范围或旧数字\\
正式交付 & 同时检查事实、工具轨迹与可预览产物 & 常出现“回答完成、产物未过硬门”\\
\bottomrule
\end{tabularx}
\subsection*{为何采用“硬门”而非主观打分}
一次运行只有同时满足下列条件才计为完成：
\[
\text{Pass}=\text{进程完成}\land\text{预算合规}\land\text{事实保留}\land\text{范围正确}\land\text{工具顺序正确}\land\text{产物可预览}.
\]
这使得“写得像一份报告”和“确实完成科研任务”被清楚地区分开。每项失败都能回到具体断言与原始事件。
\begin{darkbox}
\textbf{可直接引用的结论}\par
在固定康复科研任务、固定模型标签和统一验收口径下，Galen 的状态化、契约化执行层显著提高了任务闭环完成率，同时降低了工具调用与完成时延。
\end{darkbox}
\clearpage

\kicker{02  /  实验设计}
\section*{同任务、同输入、同指标，比较执行架构}
\begin{tabularx}{\textwidth}{@{}p{31mm}Y Y Y@{}}
\toprule
\textbf{任务组} & \textbf{PCTX 01--05} & \textbf{PDATA 01--05} & \textbf{PTOOL 01--05}\\
\midrule
考察能力 & 项目记忆、修订、范围切换、干扰抵抗 & 字段规范、缺失/重复、时间轴、来源冲突 & 检索、读写、失败恢复、正式交付\\
验收对象 & 当前事实与旧范围隔离 & 数据结构与质量记录 & 工具序列、预算、链接与产物\\
运行覆盖 & 5 卡 × 5 次 & 5 卡 × 5 次 & 5 卡 × 5 次\\
\bottomrule
\end{tabularx}
\subsection*{比较对象与覆盖}
\begin{tabularx}{\textwidth}{@{}Y r r Y@{}}
\toprule
\textbf{系统} & \textbf{任务卡} & \textbf{运行数} & \textbf{角色}\\
\midrule
Galen Full & 15 & 75 & 完整状态层、数据契约与执行契约\\
Codex CLI + DeepSeek & 15 & 75 & 完整重复的外部通用 Agent 基线\\
OpenCode + DeepSeek & 15 & 15 & 全任务单次兼容性先导\\
OpenCode + DeepSeek & 3 & 15 & 代表性任务 Repeat-5 稳定性观察\\
\bottomrule
\end{tabularx}
\subsection*{可重复性控制}
\begin{tabularx}{\textwidth}{@{}Y Y@{}}
\begin{card}{paleBlue}\textbf{冻结输入}\par 同一任务文本、fixture、必需事实、禁用事实、允许工具和交付文件。\end{card}&
\begin{card}{paleTeal}\textbf{冻结判定}\par 相同成功条件、工具预算、模型请求预算、超时与产物检查。\end{card}\\
\begin{card}{paleGold}\textbf{原始记录}\par 每次运行保留 JSON 事件、工具轨迹、输出文件、断言结果与耗时。\end{card}&
\begin{card}{white}\textbf{统计方法}\par Wilson 95\% 区间；以任务卡为重采样单位进行 20,000 次 bootstrap。\end{card}
\end{tabularx}
\small\color{muted}说明：OpenCode 为先导覆盖，未与 75 次重复结果合并；主比较为 Galen 与 Codex 两套完整重复系统。\color{ink}
\clearpage

\kicker{03  /  主结果}
\section*{Galen 在严格交付门槛下保持稳定完成}
\begin{figure}[H]
\centering
\includegraphics[width=0.99\textwidth]{D:/DEV/Galen-new/output/pdf/galen-framework-comparison/figures/fig1_success.pdf}
\caption{冻结任务卡协议下的硬门完成率。误差线为 Wilson 95\% 区间；OpenCode 的单次先导与重复子集分开展示。}
\end{figure}
\begin{tabularx}{\textwidth}{@{}Y r r r@{}}
\toprule
\textbf{系统} & \textbf{通过/总数} & \textbf{完成率} & \textbf{Wilson 95\% 区间}\\
\midrule
\textbf{Galen} & \textbf{74/75} & \textbf{98.7\%} & \textbf{92.8--99.8\%}\\
Codex + DeepSeek & 21/75 & 28.0\% & 19.1--39.0\%\\
OpenCode + DeepSeek（先导） & 9/15 & 60.0\% & 35.7--80.2\%\\
\bottomrule
\end{tabularx}
\begin{card}{paleTeal}
\textbf{配对结果}\par
相同 case/run index 下，Galen 独有通过 53 次，Codex 独有通过 0 次，共同通过 21 次，共同失败 1 次。按任务卡 bootstrap 的完成率差为 \textbf{70.7 个百分点}，95\% 区间为 \textbf{52.0--88.0 个百分点}。
\end{card}
\small\color{muted}结果解释范围：该实验比较的是固定康复科研工作流中的执行可靠性，不是对基础模型通用智能的排名。\color{ink}
\clearpage

\kicker{04  /  工具经济性}
\section*{更少的工具动作，完成更完整的交付}
\begin{figure}[H]
\centering
\includegraphics[width=0.99\textwidth]{D:/DEV/Galen-new/output/pdf/galen-framework-comparison/figures/fig2_efficiency.pdf}
\caption{相同 15 卡 × 5 次重复协议下的工具调用与端到端时延。}
\end{figure}
\begin{tabularx}{\textwidth}{@{}Y r r r@{}}
\toprule
\textbf{指标} & \textbf{Galen} & \textbf{Codex + DS} & \textbf{相对变化}\\
\midrule
平均可观察工具调用 & 2.11 & 6.12 & \textcolor{teal}{\textbf{减少 65.5\%}}\\
中位端到端时延 & 9.2 s & 19.3 s & \textcolor{teal}{\textbf{缩短 52.3\%}}\\
硬门完成率 & 98.7\% & 28.0\% & \textcolor{blue}{\textbf{提高 70.7 个百分点}}\\
\bottomrule
\end{tabularx}
\begin{card}{paleBlue}
\textbf{工具经济性为何重要}\par
科研 Agent 的调用成本不仅是 Token。无目的地遍历目录、重复检索、反复改写同一文件，会同时拉长时间、增加失败面并削弱结果可复核性。Galen 通过项目状态、工具白名单、顺序约束和调用预算，把每次读取、检索与写入绑定到明确的研究状态。
\end{card}
\begin{darkbox}
\textbf{评审判断}\par
Galen 的优势不是“少调用所以少做事”，而是在更高完成率的前提下，以更少工具动作交付了通过事实、范围、轨迹和产物四重检查的结果。
\end{darkbox}
\clearpage

\kicker{05  /  失败边界}
\section*{优势集中在哪里，薄弱点也清楚可见}
\begin{figure}[H]
\centering
\includegraphics[width=0.99\textwidth]{D:/DEV/Galen-new/output/pdf/galen-framework-comparison/figures/fig3_failure_boundary.pdf}
\caption{逐任务卡硬门完成率。Galen 与 Codex 每卡 $n=5$；OpenCode 为每卡 $n=1$ 的兼容性先导。}
\end{figure}
\begin{tabularx}{\textwidth}{@{}p{32mm}Y Y@{}}
\toprule
\textbf{能力区} & \textbf{Galen 表现} & \textbf{通用基线主要失败模式}\\
\midrule
上下文 PCTX & 绝大多数任务 5/5 & 修订后重新带回旧病种、旧样本量或旧时间窗\\
数据治理 PDATA & 五张任务卡均稳定完成 & 能解释规则，但未在预算内生成合规数据产物\\
工具与交付 PTOOL & 仅 PTOOL02 出现 1 次预算越界 & 检索循环、工具顺序不符、失败恢复链中断\\
\bottomrule
\end{tabularx}
\begin{card}{paleGold}
\textbf{保留的真实失败}\par
Galen 在 PTOOL02 有 1 次未通过：模型/工具请求超过任务设定的双预算。该记录未被重试覆盖，直接进入最终统计。它推动了长任务预算提示、细粒度进度与恢复点的后续改进。
\end{card}
\clearpage

\kicker{06  /  架构消融}
\section*{上下文能力来自状态层，而不是“模型多记一点”}
\subsection*{配对上下文实验：完整历史 vs. 结构化上下文包}
\begin{tabularx}{\textwidth}{@{}Y r r r r@{}}
\toprule
\textbf{变体} & \textbf{通过} & \textbf{旧范围失败} & \textbf{平均输入 Token} & \textbf{TTFR 中位区间}\\
\midrule
完整历史 baseline & 7/25 & 18 & 73,462 & 1.72--1.90 s\\
\textbf{Galen 上下文包} & \textbf{25/25} & \textbf{0} & \textbf{15,601} & \textbf{1.30--1.59 s}\\
\bottomrule
\end{tabularx}
\begin{card}{paleTeal}
结构化上下文包在本组实验中将输入 Token 平均值降低 \textbf{78.8\%}，同时把完成率从 \textbf{28\%} 提高到 \textbf{100\%}。关键不是简单截断历史，而是显式保留研究问题、当前范围、排除方向、数据事实和已确认决策。
\end{card}
\subsection*{关键任务 Repeat-5 消融}
\begin{tabularx}{\textwidth}{@{}Y r r r r@{}}
\toprule
\textbf{变体} & \textbf{PCTX02} & \textbf{PCTX03} & \textbf{合计} & \textbf{主要现象}\\
\midrule
\textbf{Full Galen} & \textbf{5/5} & \textbf{5/5} & \textbf{10/10} & 当前范围完整保留\\
Stateless & 1/5 & 1/5 & 2/10 & 旧样本量、旧量表回流\\
Generic same runtime & 0/5 & 1/5 & 1/10 & 上下文泄漏并触发预算超限\\
No execution contract & 4/5 & 4/5 & 8/10 & 工具调用或事实约束偶发失守\\
\bottomrule
\end{tabularx}
\begin{darkbox}
\textbf{消融结论}\par
同一运行环境中关闭状态层后，两个高难度范围切换任务的通过数由 10/10 降至 2/10；这直接说明 Galen 的连续研究能力来自可执行项目状态，而不是聊天记录堆叠。
\end{darkbox}
\clearpage

\kicker{07  /  迭代与工程验收}
\section*{测试不仅证明结果，也直接驱动版本修复}
\subsection*{长任务可靠性回归}
\begin{tabularx}{\textwidth}{@{}Y r r r@{}}
\toprule
\textbf{指标} & \textbf{修复前} & \textbf{修复后} & \textbf{变化}\\
\midrule
硬门通过 & 27/30 & 30/30 & +3 次\\
Pass$^5$ & 66.7\% & 100.0\% & +33.3 个百分点\\
Wilson 95\% 下界 & 74.4\% & 88.6\% & +14.2 个百分点\\
总耗时 P95 & 107.3 s & 21.5 s & \textcolor{teal}{\textbf{降低 80.0\%}}\\
\bottomrule
\end{tabularx}
\begin{card}{paleBlue}
\textbf{修复来源}\par
Repeat-5 回归曾捕获两次读取文件后的长时间流式空转，以及一次真实观察 ID 未被带入产物。加入流空闲看门狗、可观察重试与病例 ID 契约后，6 张任务卡 × 5 次全部通过。
\end{card}
\subsection*{当前代码质量门}
\begin{tabularx}{\textwidth}{@{}Y r@{}}
\toprule
\textbf{检查项} & \textbf{结果}\\
\midrule
Rust workspace 编译检查 & \statusok\\
Galen 核心单元测试 & \textcolor{teal}{\bfseries 219 passed / 0 failed / 1 ignored}\\
前端 TypeScript 静态检查 & \statusok\\
PDF 与产物可预览链路 & \statusok\\
\bottomrule
\end{tabularx}
\begin{card}{paleGold}
\textbf{版本演进证据}\par
严格对齐批次中，Galen 曾取得 72/75（96.0\%）；修复工具预算与交付链后，最终冻结批次达到 74/75（98.7\%）。报告同时保留两批记录，避免只展示最佳一次运行。
\end{card}
\clearpage

\kicker{08  /  复现与材料索引}
\section*{每一个数字都能回到原始记录}
\begin{tabularx}{\textwidth}{@{}p{39mm}Y@{}}
\toprule
\textbf{证据层} & \textbf{仓库位置}\\
\midrule
一键核验入口 & \code{07-其他材料/run\_judge\_benchmark.ps1}\\
三张代表案例 & \code{07-其他材料/evidence-data/representative-cases/}\\
冻结任务卡与 fixture & \code{evals/cases/paired/}\\
Galen 最终 75 次记录 & \code{evals/runs/large-full-galen-20260912-final.jsonl}\\
Codex 最终 75 次记录 & \code{evals/runs/large-codex-deepseek-flash-20260912-complete.jsonl}\\
最终统计摘要 & \code{evals/runs/large-comparison-20260912-final.md}\\
上下文配对实验 & \code{evals/runs/paired-v1/20260911-formal-v1/}\\
架构消融记录 & \code{evals/runs/architecture-v1/}\\
长任务 Repeat-5 回归 & \code{evals/reports/}\newline\code{ais-t2-source-closed-watchdog-pro-repeat5-final-2026-08-25.md}\\
统计图与 LaTeX 报告 & \code{output/pdf/galen-framework-comparison/}\\
\bottomrule
\end{tabularx}
\subsection*{随本目录提供的评审核验材料}
\begin{itemize}
  \item \code{run\_judge\_benchmark.ps1}：无须 API Key，重算 150 次记录的完成率、工具调用与时延；
  \item \code{evidence-data/}：两份完整 JSONL、三张代表任务卡、双方轨迹、Galen 实际产物与 SHA-256 清单；
  \item \code{证据总表.md}：从每项申报主张直接跳转到对应证据和核验动作；
  \item \code{paired-benchmark/}：15 张冻结任务卡、固定输入与自动判定条件；
  \item \code{run\_demo\_checks.ps1}：基础编译、单元测试与前端检查入口。
\end{itemize}
\begin{darkbox}
\textbf{最终结论}\par
Galen 已经形成“固定任务—自动执行—硬门判定—失败定位—版本修复—重复验证”的完整实验闭环。多智能体对比说明了专用科研执行架构的价值，消融实验说明了优势来源，原始记录与脚本保证结论可以复核。
\end{darkbox}
\vfill
\begin{center}
\href{https://github.com/Drehabwen/Galen/tree/galen-research-workbench}{\color{blue}\bfseries 查看 Galen 公开仓库与提交分支}\par
\small\color{muted}Galen Research Workbench  /  2026
\end{center}
\end{document}
